\par\Needspace{6\baselineskip}\medskip
\SourceLeaf{157}\pdfbookmark[2]{112.}{article-112}
\begin{center}112.\end{center}
Es seien $x''$, $y''$ und $z''$ die Koordinaten eines dritten Ortes, $u''$ sein Breitenargument und $r''$ sein Radiusvektor. Wir bezeichnen die Größen
\[
r'r''\sin(u''-u'),\qquad rr''\sin(u''-u),\qquad rr'\sin(u'-u),
\]
die die doppelten Flächeninhalte der Dreiecke zwischen dem zweiten und dritten, dem ersten und dritten sowie dem ersten und zweiten Radiusvektor sind, mit $n$, $n'$ beziehungsweise $n''$. Für $x''$, $y''$ und $z''$ erhält man somit Ausdrücke, die den im vorigen Artikel für $x$, $y$, $z$ und für $x'$, $y'$, $z'$ angegebenen ähnlich sind. Mithilfe von Lemma I in Artikel 78 lassen sich daraus leicht die folgenden Gleichungen ableiten.
\[
\begin{aligned}
0&=nx-n'x'+n''x''\\
0&=ny-n'y'+n''y''\\
0&=nz-n'z'+n''z''.
\end{aligned}
\]
Es seien nun $\alpha$, $\alpha'$ und $\alpha''$ die geozentrischen Längen des Himmelskörpers an den drei Orten, $\beta$, $\beta'$ und $\beta''$ seine geozentrischen Breiten und $\delta$, $\delta'$ und $\delta''$ seine auf die Ekliptik projizierten Abstände von der Erde. Ferner seien $L$, $L'$ und $L''$ die entsprechenden heliozentrischen Längen der Erde und $B$, $B'$ und $B''$ ihre Breiten. Diese Breiten setzen wir nicht $=0$, damit sich die Parallaxe berücksichtigen und nach Belieben anstelle der Ekliptik jede andere Ebene wählen lässt. Schließlich seien $D$, $D'$ und $D''$ die auf die Ekliptik projizierten Abstände der Erde von der Sonne. Drückt man nun $x$, $y$ und $z$ durch $L$, $B$, $D$, $\alpha$, $\beta$ und $\delta$ aus und verfährt ebenso mit den Koordinaten des zweiten und dritten Ortes, so nehmen die vorangehenden Gleichungen die folgende Form an.
\[
\begin{aligned}
1)\quad 0&=n(\delta\cos\alpha+D\cos L)-n'(\delta'\cos\alpha'+D'\cos L')\\
&\qquad+n''(\delta''\cos\alpha''+D''\cos L'')\\[.5ex]
2)\quad 0&=n(\delta\sin\alpha+D\sin L)-n'(\delta'\sin\alpha'+D'\sin L')\\
&\qquad+n''(\delta''\sin\alpha''+D''\sin L'')\\[.5ex]
3)\quad 0&=n(\delta\operatorname{tang}\beta+D\operatorname{tang}B)-n'(\delta'\operatorname{tang}\beta'+D'\operatorname{tang}B')\\
&\qquad+n''(\delta''\operatorname{tang}\beta''+D''\operatorname{tang}B'').
\end{aligned}
\]
Betrachtet man hier $\alpha$, $\beta$, $D$, $L$ und $B$ sowie die entsprechenden Größen für die beiden übrigen Orte als bekannt und teilt die Gleichungen durch $n$, $n'$ oder $n''$, so bleiben fünf Unbekannte. Von diesen kann man zwei eliminieren oder aus zwei beliebigen die übrigen drei bestimmen. Auf diese Weise führen die drei Gleichungen zu vielen wichtigen Folgerungen, von denen wir hier einige besonders bemerkenswerte entwickeln.
